\capitulo{7}{Conclusiones y Líneas de trabajo futuras}

En este último capítulo se exponen las conclusiones generales derivadas del desarrollo del sistema HIVES, analizando el grado de consecución de los objetivos planteados y evaluando críticamente los retos de ingeniería de software superados durante la integración. Asimismo, se definen las líneas de investigación y desarrollo futuras orientadas a la evolución tecnológica del dispositivo y del software predictivo.

\section{Conclusiones}

Tras la finalización del proyecto, la conclusión principal es que se han cumplido satisfactoriamente la totalidad de los objetivos técnicos y funcionales establecidos en la especificación inicial del sistema. Desde una perspectiva académica y profesional, el desarrollo del analizador HIVES ha constituido un hito formativo de gran valor, al haber permitido consolidar y poner en práctica de manera unificada en un único proyecto tangible los conceptos transversales adquiridos a lo largo del Grado en Ingeniería Informática (tales como la arquitectura de software, la gestión de bases de datos, el desarrollo de interfaces gráficas reactivas y la aplicación práctica de la inteligencia artificial).

En lo relativo al comportamiento operativo y al rendimiento de la solución, el sistema ha demostrado una alta fiabilidad en la adquisición y procesamiento de datos. No obstante, la integración de la capa de comunicación hardware-software supuso uno de los mayores retos del proyecto. En las fases iniciales de pruebas, surgieron problemas de sincronización en la transmisión de flujos serie desde el microcontrolador hacia la estación de trabajo. El principal obstáculo detectado se centró en escenarios de desconexiones intermitentes debidos a conflictos en el acceso concurrente al puerto serie virtual. Durante el flujo de desarrollo, la monitorización simultánea mediante el monitor serie del \ruta{Arduino IDE} y la escucha activa del controlador implementado en \ruta{Visual Studio Code} provocaban la saturación y el bloqueo del buffer serie del hardware. Este inconveniente técnico se solventó con éxito mediante la depuración estricta del protocolo de comunicación y la correcta delegación de tareas asíncronas a través de hilos de ejecución independientes, garantizando la resiliencia de la interfaz gráfica y la estabilidad del canal serie en la versión final de producción.

Finalmente, cabe destacar el impacto formativo derivado de las limitaciones operativas del conjunto de datos. La necesidad de explorar y verificar el comportamiento de redes neuronales artificiales densas con un volumen de 42 muestras impulsó la investigación activa en el ámbito de la computación y el entrenamiento en la nube. La utilización de entornos cloud avanzados como \ruta{Google Colab} permitió experimentar con el rendimiento algorítmico de forma ágil y deslocalizada. Este campo técnico, que habitualmente cuenta con una exposición limitada dentro de los planes de estudio de la titulación, ha resultado ser un área de especialización sumamente interesante, aportando conocimientos críticos sobre las  infraestructuras para proyectos complejos de computación intensiva.

\section{Líneas de trabajo futuras}

A partir de los resultados y de la experiencia práctica obtenida en este Trabajo de Fin de Grado, se proponen las siguientes líneas de trabajo futuras para expandir las capacidades funcionales y tecnológicas del sistema:

\begin{itemize}
    \item \textbf{Ampliación del conjunto de datos empíricos y revalidación del clasificador:} El objetivo prioritario e inmediato consiste en expandir de forma sustancial el volumen del \textit{dataset} de firmas espectrales reales obtenidas directamente con el sensor AS7265x. Tal y como se analiza en detalle en la Sección~5.2 de la memoria, las métricas de clasificación reportadas en el presente
    trabajo fueron obtenidas sobre un conjunto de datos mayoritariamente sintético ---generado mediante difusión gaussiana a partir de únicamente 42 muestras reales---, por lo que su validez  en condiciones  reales es limitada. La incorporación continua de nuevas muestras certificadas mediante análisis melisopalinológico independiente permitirá, en primer lugar, reentrenar el clasificador sobre datos  suficientes y, en segundo lugar, disponer de un conjunto de prueba independiente con el que evaluar el rendimiento real del sistema. Solo a partir de ese punto podrá considerarse que las métricas de exactitud del clasificador constituyen evidencia  válida para un contexto de aplicación práctica o comercial.
    
    \item \textbf{Evolución del hardware espectrométrico (Bandas extremas UV e IR):} El hardware actual (AS7265x) ha demostrado una excelente capacidad para extraer características y predecir parámetros integrados en el modelo de \textit{Machine Learning} dentro del rango de 410~nm a 940~nm. No obstante, estudios académicos consolidados~\cite{sciencedirect_espectro_miel} demuestran que el análisis de muestras de miel en longitudes de onda situadas por debajo del espectro visible (ultravioleta profundo) y muy por encima de este (infrarrojo medio y cercano, NIRS~\cite{guillen2017nirs}) revela información altamente discriminativa sobre compuestos fenólicos y origen floral. Por ello, una evolución natural del prototipo sería la integración de sensores complementarios que abarquen estas bandas extremas.
    
    \item \textbf{Migración hacia un ecosistema móvil y portátil:} Con el fin de maximizar la usabilidad en entornos de campo (como en explotaciones apícolas productivas), se contempla la sustitución de la arquitectura de escritorio local por un despliegue portátil. Esta evolución requerirá el desarrollo de una aplicación móvil multiplataforma y la reingeniería del módulo de comunicaciones para reemplazar la conexión física USB por un protocolo inalámbrico mediante tecnología \ruta{Bluetooth}, dotando al hardware de una autonomía completa.
\end{itemize}
